(1) 驻点由 \(f_P=a(T-T_c)P+bP^3-E=0\) 给出。取 \(E=0\)，根为 \(P=0\) 及在 \(T<T_c\) 的 \(P=\pm\sqrt{-a(T-T_c)/b}\)。曲率 \(f_{PP}=a(T-T_c)+3bP^2\) 必须为正才能局部稳定，因此 \(T>T_c\) 只有 \(P=0\) 稳定；在 \(T<T_c\)，零极化变为极大值，两非零根为简并最小值。\(\boxed{|P_0|\propto(T_c-T)^{1/2}}\)，这是平均场临界指数 \(1/2\)，并不是任意维数实材料的普适精确指数。

(2) 对平衡方程隐式微分得
\((a(T-T_c)+3bP^2)\,\partial P/\partial E=1\)，所以稳定支的响应
\(\chi=1/[a(T-T_c)+3bP^2]\)。在 \(T>T_c,E=0\)，\(\chi_+=1/[a(T-T_c)]\)；在 \(T<T_c\) 的选定非零支，代入 \(bP_0^2=-a(T-T_c)\)，得
\[
\boxed{\chi_-=\frac1{2a(T_c-T)}}.
\]
两侧均按 \(|T-T_c|^{-1}\) 发散，但同等温差时高温幅度是低温的两倍；响应定义必须在选定稳定相上求导，而非对跨跃迁的混合平均直接求导。

(3) 旋节点满足 \(f_P=f_{PP}=0\)。令 \(A=a(T_c-T)>0\)，则
\(P_{\rm sp}=\pm\sqrt{A/(3b)}\)，对应场
\(E_{\rm sp}=-AP_{\rm sp}+bP_{\rm sp}^3=-(2A/3)P_{\rm sp}\)，故
\[
\boxed{|E_{\rm sp}|=\frac{2A^{3/2}}{3\sqrt{3b}}}.
\]
这是均匀单畴自由能的局部极小消失极限；实际切换可在较小场通过成核、缺陷与畴壁运动实现，故不得将理想旋节场无条件当作实验矫顽场。

(4) 低温零场在 \(P=0\) 的自由能取0，在 \(\pm P_0\) 的值 \(f(P_0)=-A^2/(4b)\)，因此能垒为 \(\Delta f=A^2/(4b)\)。选定正负态施加微弱电场后自由能差主阶为 \(\Delta f_{+-}\simeq-2EP_0\)，外场选择与自身方向同号的极化。若 \(b\to0^+\) 而保留负二次系数，自发极化与该四次模型的能垒形式发散，说明更高偶次项和材料饱和机制不可忽略，不能把四阶Landau截断继续当作定量模型。
